\section{难点，以及我是怎么改的}

这一节按"我本来以为是什么样 → 哪里不对 → 改了什么 → 现在是什么数"来写。选的都是真卡住过我的地方。

\subsection{玩家不打开聊天，就等于没有教练}

第一版做出来我自己试了两局，发现一件事：小芽的所有能力都在聊天面板里，而我全程没打开过那个面板。一个要玩家主动召唤才存在的助手，在真实游玩里等于不存在。

改法是把主入口从聊天框挪到游戏事件上。现在的规则是：开局和关键风险变化时，在战斗页底部出现一行短提示，解释收在二级入口里按需展开；玩家出招之后，旧提示立刻隐藏，等新局面重新评估。聊天框降级成补充入口。

这条改动有个副作用值得说：提示一多就变成噪音。所以门控比内容更难做，见后面第五个难点。

\subsection{数字不能让模型算}

我一开始让模型看局面然后给建议，试出来的回答里数字是错的。不是偶尔错，是稳定地错——它会把不存在的技能名编出来，会把伤害算错，会在两个都能击倒的技能里挑面板数字大的那个。

后来把职责切开：引擎枚举双方所有合法行动，对每种组合同跑一遍共享结算器，得到平均收益和最坏分支。模型拿到的是一份事实包，里面只有引擎算出来的数字和每张战术卡的反例，它负责把这些讲成取舍。

有个例子能说明为什么这个分工是必须的。同一个残血目标，火花直接伤害 52，追猎 75，两个都够击倒，都耗 2 豆。只看面板会选追猎。但追猎的额外伤害来自"目标已灼烧"这个前提，对方吃药、换宠、防御都会让前提失效。这种判断交给模型是猜，交给枚举是算。

\subsection{检索没有赢过关键词}

我做过一个假设：把检索从关键词换成"关键词+多语言向量"的混合方案，口语化提问的命中率会提高。为了验证它，我写了一个对照实验，同一批查询分别跑无检索、纯关键词、纯向量、混合四档。

结果是混合方案没有赢。测试集 14 条正例上，关键词命中 11 条（Hit@3 = 0.786），混合也是 11 条（0.714），差异来自排序而不是命中的数量。负例拒答上两者更接近，关键词拒掉 1 条，纯向量拒掉 2 条。

这里我差点犯一个错。看到混合没赢，我的第一反应是"检索没用，RAG 是伪需求"。但把四档拆开看就发现不是这么回事：纯向量在负例拒答上比关键词好，混合的 MRR 从 0.667 升到 0.750，排序确实更好了。问题是这些改善都发生在 14 条正例、2 条负例的样本上。一条查询的差别就能让 Hit@3 动 7 个百分点。

所以我没有拿这个结果去证明任何事。既没有换成向量检索，也没有宣称 RAG 无效，而是把它记成"这次测量没测出差别"，并把关键词保留为默认。后来读 BEIR 那篇才知道这类现象在检索领域很常见：BM25 在跨域零样本场景下本来就是很难打败的基线，稠密和稀疏检索模型在原论文里的表述是 "often underperform"。

真正的教训是我一开始把评测集做得太小。20 条查询、4 条开发、16 条测试，这个规模连自己的结论都支撑不住。

\subsection{模型会说过期的话}

这个问题的表现是：玩家已经出招了，模型那边的解释才回来，界面上出现一句针对已经过去的局面的建议。

我试过在提示词里写"如果局面已经变化就不要回答"，没用。模型看不到自己那句回答会在什么时候被显示出来，这个判断只能由程序做。

改法是把每个请求绑定到具体的局面版本：请求发出时记下对局 ID、回合号、规则版本和上下文 epoch，返回时逐项比对。任何一项变了就丢弃结果，不显示、也不播报。这条对语音同样适用。

后来发现还有一个更隐蔽的变体：不是结果过期，是我自己的版本号没同步。服务端校验上下文模式的白名单是 \texttt{camp / pve / pvp-live}，我加本地对战时客户端开始发 \texttt{pvp-local}，但白名单没跟着改。结果是本地对战里的模型请求全部被 400 拒绝、静默回退成本地规则。规则建议照常工作，界面没有任何异常提示，这个模式从头到尾没有产生过一次模型解释。它是靠逐条核对才发现的，现在有一个回归测试遍历客户端可能发送的每一种模式。

\subsection{模型会编道具名，也会把取消的行动说成命中}

有两类错误是数字校验抓不到的。

第一类是道具名。这个游戏里只有回复药、净化药、能量果三样。模型会把净化药说成"解药"、能量果说成"以太"。数字全对，引用全对，只有名字是错的——因为它没有数字可校验。

第二类是因果。宠物倒下时，它这一回合已经排好的行动会被取消。模型会把这次取消描述成"命中了"。同样没有数字可比。

这两类都靠加检查解决：一是维护一份别名黑名单，命中就判不合格；二是读回合事件，如果某个行动被记为取消，那么正文里就不能出现这个行动造成了伤害。判不合格之后走既有的降级路径，显示规则结论而不是错误描述。

\subsection{什么时候不说话，比说什么更难}

提示做出来之后，问题从"说得对不对"变成"该不该说"。我给的门控是逐项判断的：玩家是否在前台、是否已经出招、是否已经忽略过、这个回合是否已经提示过、上次提示距今多久、当前风险是否构成可行动的风险。安静档和"本局不再提醒"永远优先，不被模型推断覆盖。

还有一条是我从一篇教育数据挖掘的论文里学到的。那篇讲 assistance dilemma，作者报告说他们早期版本的预测器没有把"学生用了提示"这件事算进去，结果因为提示本身会把人推向高效步骤，用过的学生反而被判定成"不需要帮助"。他们的修正是把两者分开。

小芽里对应的是同一件事：判断玩家是否掌握某个知识点时，只用他没有被提示过的时候做出的选择。代码里就是一个过滤条件 \texttt{!e.prompted}。我是先按直觉这么设计的，后来才发现这个踩坑在文献里有记录。

\subsection{对局里的对手，我没有交给模型}

题目场景是 PVP，我一开始的想法是让模型来当对手，理由是"这样显得 AI 更强"。查了资料之后放弃了这个想法。

放弃的理由不是我懒，是证据方向相反。PokéLLMon 是一个纯靠语言模型打宝可梦的工作，阶梯赛胜率 49\%。PokéChamp 换成"LLM 辅助 minimax 搜索"之后，对最强规则型 bot 胜率是 84\%，而且不需要任何额外训练。Metamon 进一步用离线强化学习加 Transformer，论文里说它的表现超过了 LLM 智能体和强启发式搜索引擎。也就是说在这个游戏类型里，决定胜负的是搜索和训练，不是语言能力。

所以对手出招我用引擎的启发式：轻松档随机、标准档优先伤害和治疗、挑战档枚举双方合法行动并比较一回合结果。模型完全不参与。这样做的另一个好处是本地出招零延迟、零成本，也不会返回非法行动。

顺带说一个由此产生的设计问题。做了本地对战之后我发现，"电脑对手"和"训练模式"在玩法上其实是同一个东西——同一套 AI、同一套结算，唯一的区别是教练闭不闭嘴。这个区别是真的，但我原来把它藏在名字后面，第一次给用户看的时候被直接问了"这两个有啥区别"。现在两个入口的名字直接写明差别：训练是"按关卡挑战"，对局是"对手自由配队"。

\subsection{教练在 PVP 里该不该说话，我改错了两次}

这一条是我这次最值得记的。

题目对军师的描述是"在对战/阵容场景给出决策建议"，而场景设定写的是 PVP 对战。也就是说军师本来就该在对战里工作。

我第一版的做法是：只要是对局模式，教练一律不介入。理由是"两人共用一块屏，给一方提示就是不公平"。这个理由听起来成立，还把"公平性设计"写进了文档当成亮点。

问题是这样等于把主功能在题目指定的场景里禁用了。改法是把判据从"是不是对局模式"换成"有没有第三方受影响"：对手是真人时闭麦，对手是电脑时照常工作。

改完之后又有人指出这个版本也不自洽：电脑对手存在的意义就是模拟真人，那它和真人局的行为就该一致，否则这个模拟本身就是假的。这句话是对的。最终改成本地对战一律允许教练，只保留线上竞技的闭麦开关。真正需要守住的底线不是"整局不说话"，而是"不读取对方待执行的动作"——那条项目一直没破，也有测试。

前后改了三次，前两次都是我把自己的推论当成了结论。

\subsection{"公平"这件事最后是靠对称解决的，不是靠闭麦}

接着上一条。既然本地对战两边都能用教练，那公平性就不需要靠禁用某个功能来实现了。

具体做法是双方各有一条教练条，跑同一套引擎、同一套规则，各自只分析自己那一侧。对面用的是把 player 和 enemy 对调的镜像上下文，所以它看到的是对方的局面，也读不到任何隐藏信息——待执行的动作本来就不在上下文里。

实测两条教练在同一回合给的结论是不同的，而且各自正确：一方建议"火花"，另一方建议"换上晶角犀"。

\subsection{对手要像真人，就得看不见我在干什么}

本地对战的对手如果是 AI，它和真人的差别会体现在一个地方：真人不可能知道你选了什么。

我把这件事做成顺序上的保证，而不是提示词里的约束。每回合开始时，先让对手独立做决定——它读的是回合前的公开局面——然后才开放我方的操作面板。它决定的时刻早于我的选择，所以不可能参考我的选择。对方面板会直接写明"已独立出招（看不到你的选择）"，按钮也是禁用的，不是点了没反应。

\subsection{分屏之后，两个面板必须一模一样}

第一版分屏我只给我方面板做了标签页（技能/换宠/物品/逃跑），对方面板是一个平铺的列表。自己看着别扭，被指出之后才发现更实际的问题：右边有标题行、左边没有，导致右侧整块内容比左侧低了大约 30 像素，标签、按钮、教练条全部错位。

改成两侧走同一个渲染函数，结构完全一致，只有"逃跑"和"认输"的区别——引擎里 escape 的语义属于本地撤退，不能给对手用。中间加了一道竖线分隔。改完实测两侧的标题、标签行、技能区、教练条纵坐标完全对齐。
